\section{Listening}

\subsection{Unit 1}
\subsubsection{Conversation}
\begin{enumerate}
	\item Why does Stacey come to Dr. Pearl's office?
	\begin{choices}
		\item To get permission for sick leave.
		\item \textbf{To get permission to quit his class.}
		\item To get permission to attend his class.
		\item To get permission to get her essay back.
	\end{choices}
	
	\item What is worrying Stacey about her studies?
	\begin{choices}
		\item She cannot take a film class next semester.
		\item She cannot cover her humanities requirements.
		\item \textbf{She feels that the writing course is too challenging.}
		\item She feels that passing an engineering course is difficult.
	\end{choices}
	
	
	\item What does Dr. Pearl suggest Stacey do?
	\begin{choices}
		\item \textbf{Sign up for free tutoring in writing.}
		\item Work with him at his office every day.
		\item Go to the University Writing Center each Friday.
		\item Work with him at The Found Librarian everyday.
	\end{choices}

	\item What is Dr. Pearl's attitude toward Stacey?
	\begin{choices}
		\item \textbf{Patient.}
		\item Satisfied.
		\item Indifferent.
		\item Disappointed.
	\end{choices}

\end{enumerate}


\subsubsection{Passage}

\begin{enumerate}
	
	\item What can we learn about confidence from the passage?
	\begin{choices}
		\item \textbf{It keeps us striving for success.}
		\item It is just an attitude toward life.
		\item It needs justification from other people.
		\item It can help us find our sense of purpose.
	\end{choices}
	
	\item How can we develop confidence, according to the passage?
	\begin{choices}
		\item \textbf{By doing the things we fear.}
		\item By learning from our failures.
		\item By knowing more about ourselves.
		\item By avoiding the feeling of insecurity.
	\end{choices}
	
	\item What should we do if our life is not what we long for?
	\begin{choices}
		\item Avoid making hasty decisions.
		\item \textbf{Make changes on a daily basis.}
		\item Reassess our goals and strategies.
		\item Be patient with life and see what happens.
	\end{choices}
	
\end{enumerate}




\subsubsection{Lectures 1}

\begin{enumerate}
	DCA
	1. What did the speaker discover when teaching math in a public school?
	2. What topic interested the speaker as a psychologist?
	3. What does the speaker say about grit?
	
	1.
	
	A
	Some students could not finish homework assignments.
	
	B
	Seventh-grade math was very difficult for most students.
	
	C
	Every student could work hard to understand the material.
	
	D
	IQ wasn't the only factor that influenced students' performance.
	
	2.
	
	A
	Challenges and mental health.
	
	B
	IQ and academic performance.
	
	C
	Factors contributing to success.
	
	D
	Mental health and environment.
	
	3.
	
	A
	We have little knowledge about how to build it.
	
	B
	Grit means living life passionately like it's a sprint.
	
	C
	It's better to develop grit in students at a younger age.
	
	D
	Students' grit can improve as their learning ability grows.
	
\end{enumerate}

\subsubsection{Lectures 2}

\begin{enumerate}
	BDAA
	1. What can we learn about procrastination from the lecture?
	2. What is the benefit of breaking down tasks into smaller steps?
	3. How can we enhance our concentration?
	4. What can we learn about the Pomodoro Technique from the lecture?
	1.
	
	A
	It traps us in a cycle of regrets.
	
	B
	It weakens our motivation and efficiency.
	
	C
	It limits our overall creativity in the long run.
	
	D
	It affects our emotional health in the long run.
	
	2.
	
	A
	It can reduce the stress of our work.
	
	B
	It can help us set clear and specific goals.
	
	C
	It enables us to prioritize the achievable tasks.
	
	D
	It allows us to focus on the most important tasks first.
	
	3.
	
	A
	By surrounding ourselves with visual cues.
	
	B
	By setting both short-term and long-term goals.
	
	C
	By building a comfortable and cheerful atmosphere.
	
	D
	By playing motivational and inspirational background music.
	
	4.
	
	A
	It is an effective technique for time management.
	
	B
	It can enhance productivity by reinforcing our goals.
	
	C
	It adds a sense of urgency by increasing stress levels.
	
	D
	It is a strategy used to improve our sense of accomplishment.
\end{enumerate}


\newpage
\subsection{Unit 2}

\subsubsection{Conversation}

\begin{enumerate}
	ACCB
	1. Why is Mount Rainier so special to the man?
	2. How often did the man and his father go hiking and camping on Mount Rainier every year?
	3. What can we learn about the man's father from the conversation?
	4. What is the relationship between the two speakers?
	1.
	
	A
	It was his father's favorite mountain.
	
	B
	It is a good place for hiking and camping.
	
	C
	He loves the yellow and blue flowers there.
	
	D
	He and his friends had many happymemories there.
	
	2.
	
	A
	Once or twice.
	
	B
	Two or three times.
	
	C
	Three or four times.
	
	D
	Four or five times.
	
	3.
	
	A
	He died at the age of 80.
	
	B
	He had a very tough life.
	
	C
	He died of a heart attack.
	
	D
	He enjoyed cooking fish.
	
	4.
	
	A
	Husband and wife.
	
	B
	Boss and employee.
	
	C
	Teacher and student.
	
	D
	Boyfriend and girlfriend.
\end{enumerate}


\subsubsection{Passage}

\begin{enumerate}
	BACC
	1. What will happen if we always think we must do something in a certain way?
	2. How can we make a large project easier, according to the passage?
	3. What is the benefit of using our beds for sleeping only?
	4. What is the passage mainly about?
	1.
	
	A
	We will become less flexible.
	
	B
	We will experience more stress.
	
	C
	We will become narrow-minded.
	
	D
	We will lose our ability to think critically.
	
	2.
	
	A
	By dividing the project into smaller parts.
	
	B
	By focusing on manageable parts of the project.
	
	C
	By tackling the most difficult part of the project first.
	
	D
	By taking more time to finish the project step by step.
	
	3.
	
	A
	It improves our memory.
	
	B
	It helps us concentrate better.
	
	C
	It promotes faster and better sleep.
	
	D
	It helps develop good sleeping habits.
	
	4.
	
	A
	Tips for enjoying a peaceful life.
	
	B
	Tips for developing positive thinking.
	
	C
	Tips for reducing stress levels at school.
	
	D
	Tips for improving academic performance.
\end{enumerate}

\subsubsection{Lectures 1}

\begin{enumerate}
	BCD
	1. What is the first step toward managing jealousy?
	2. How can we alleviate jealousy in a relationship?
	3. Why has social media turned into a breeding ground for jealousy?
	1.
	
	A
	Cultivating self-esteem.
	
	B
	Acknowledging its presence.
	
	C
	Focusing on personal growth
	
	D
	Stopping making comparisons.
	
	2.
	
	A
	By building trust with our partner.
	
	B
	By figuring out the reasons behind our jealousy.
	
	C
	By expressing concerns and listening to our partner.
	
	D
	By appreciating our own unique qualities and strengths.
	
	3.
	
	A
	It leaves us with little time to focus on self-growth.
	
	B
	It discourages us from communicating with others.
	
	C
	It leads to a constant need for acceptance and recognition.
	
	D
	It leads to misguided perceptions and unhealthy comparisons.
\end{enumerate}


\subsubsection{Lectures 2}

\begin{enumerate}
	DBAD
	1. What can we learn about negative emotions from the passage?
	2. What is the result of avoiding or suppressing negative emotions?
	3. When is a suitable time to think about the causes of our negative emotions?
	4. Why does dealing with negative emotions need patience?
	
	
	1.
	
	A
	They generally fade away slowly.
	
	B
	They form a large part of our lives.
	
	C
	They are mainly associated with failures.
	
	D
	They often mark the path to great achievements.
	
	2.
	
	A
	It reduces negative emotions temporarily.
	
	B
	It prolongs the impact of negative emotions.
	
	C
	It diminishes our awareness of negative emotions.
	
	D
	It helps us better deal with negative emotions later.
	
	3.
	
	A
	When we are calm enough.
	
	B
	When we recognize their presence.
	
	C
	When we can get support from professionals.
	
	D
	When we are worrying about unnecessary things.
	
	4.
	
	A
	Negative emotions are constantly changing.
	
	B
	It takes time for negative emotions to fade away.
	
	C
	Identifying the causes of negative emotions can be difficult.
	
	D
	It takes time to find an appropriate way to manage negative emotions.
\end{enumerate}



\newpage
\subsection{Unit 3}

\subsubsection{Conversation}

\begin{enumerate}
	DABD
	1. What kind of person does the woman want for a flatmate?
	2. What do we know about the woman's attitude to cleanliness?
	3. Why does the woman think money is a problem when looking for a flatmate?
	4. What will the woman do to make sure her flatmate can afford the rent?
	1.
	
	A
	Very considerate and caring.
	
	B
	Somewhat reserved and quiet.
	
	C
	Pretty easy-going and straightforward.
	
	D
	Sociable but also aware of personal space.
	
	2.
	
	A
	She cares a lot about being neat and tidy.
	
	B
	She doesn't care whether her flatmate is tidy or not.
	
	C
	She hopes her flatmate will offer to clean up their room.
	
	D
	She thinks she should share the cleaning with her flatmate.
	
	3.
	
	A
	She cannot afford to live alone.
	
	B
	Her previous flatmate often delayed payments.
	
	C
	She needs the rent to be paid as soon as possible.
	
	D
	She thinks reaching a deal regarding rent is difficult.
	
	4.
	
	A
	Ask them about their work.
	
	B
	Ask them how much they earn.
	
	C
	Ask them to pay the rent each month on time.
	
	D
	Ask them whether they can pay three months' rent in advance.
\end{enumerate}


\subsubsection{Passage}

\begin{enumerate}
	BCC
	1. How can we build trust and understanding with our neighbors?
	2. What should we do if we are likely to cause problems for our neighbors?
	3. How should we react if our neighbors are bothering us?
	1.
	
	A
	By visiting them frequently.
	
	B
	By getting to know each other.
	
	C
	By organizing parties together.
	
	D
	By taking family vacations together.
	
	2.
	
	A
	Address problems appropriately when they occur.
	
	B
	Apologize to our neighbors when problems occur.
	
	C
	Take action beforehand to avoid potential problems.
	
	D
	Immediately cease activities that may cause problems.
	
	3.
	
	A
	We should involve the police for help.
	
	B
	We should wait patiently for their explanations.
	
	C
	We should express our concerns and discuss solutions together.
	
	D
	We should involve more family members to address the problem.
\end{enumerate}

\subsubsection{Lectures 1}

\begin{enumerate}
	CDA
	1. Who are the most likely loyal customers of a sandwich shop, according to the lecture?
	2. Which of the following methods is crucial to building a genuine connection between a business and the local community?
	3. Which of the following strategies is effective in reinforcing the connection between a business and the local community?
	1.
	
	A
	People who love sandwiches.
	
	B
	People who are offered coupons.
	
	C
	Local residents living down the street.
	
	D
	Office workers living in the community.
	
	2.
	
	A
	Offering coupons.
	
	B
	Distributing free lunches.
	
	C
	Handing out free samples.
	
	D
	Sponsoring local organizations.
	
	3.
	
	A
	Creating loyalty programs.
	
	B
	Offering more discounts to customers.
	
	C
	Choosing several advertising platforms.
	
	D
	Working with the most suitable marketing company.
\end{enumerate}


\subsubsection{Lectures 2}

\begin{enumerate}
	ADCB
	1. Why are well-defined boundaries beneficial?
	2. What does the speaker say about boundaries in the past?
	3. What is the possible consequence of overemphasizing boundaries in today’s world?
	4. How should we address boundaries in modern times, according to the lecture?
	1.
	
	A
	They help protect privacy and prevent trespassing.
	
	B
	They separate one's property from their neighbors'.
	
	C
	They prevent the occurrence of neighborhood disputes.
	
	D
	They enhance safety by defining the zones for specific activities.
	
	2.
	
	A
	They were solid and well-defined.
	
	B
	They were more visible and tangible.
	
	C
	They were not reliable in terms of safety.
	
	D
	They were not harmful to neighborly communication.
	
	3.
	
	A
	There will be fewer casual interactions.
	
	B
	There will be fewer gatherings and parties.
	
	C
	There will be fewer free exchanges of ideas.
	
	D
	There will be a misleading sense of security.
	
	4.
	
	A
	By making more efforts to prevent invisible boundaries.
	
	B
	By overcoming the challenges posed by invisible boundaries.
	
	C
	By striking a balance between invisible and physical boundaries.
	
	D
	By removing all forms of boundaries to promote communication.
\end{enumerate}

\newpage
\subsection{Unit 4}

\subsubsection{Conversation}

\begin{enumerate}
	CBCB
	1. What is the man's problem?
	2. How does the woman usually find inspiration for her new paintings?
	3. What does the man think of the woman's way of finding inspiration?
	4. What does the woman advise the man to do at the train station?
	1.
	
	A
	He has difficulty writing personal stories.
	
	B
	He thinks the creative writing class is boring.
	
	C
	He doesn't know what to write for his fictional writing assignment.
	
	D
	He finds it challenging to write three creative stories in two months.
	
	2.
	
	A
	She goes to the train station.
	
	B
	She goes for long walks in nature.
	
	C
	She engages in conversations with artists.
	
	D
	She observes dramatic goodbyes and romantic reunions.
	
	3.
	
	A
	It's a useless old trick.
	
	B
	It's a pretty good method.
	
	C
	It might not work for him.
	
	D
	It only works for artists and writers.
	
	4.
	
	A
	Listen to others' conversations.
	
	B
	Sit in the lobby and watch others.
	
	C
	Talk with strangers to learn their stories.
	
	D
	Share the plot of his story with passers-by.
\end{enumerate}


\subsubsection{Passage}

\begin{enumerate}
	DDB
	1. What do people disagree about regarding scientific research?
	2. Why should we value research with unknown benefits?
	3. What does the speaker say about computer research?
	1.
	
	A
	Whether research should be conducted for its immediate benefits.
	
	B
	Whether research should continue even without financial support.
	
	C
	Whether the government should distribute funds for research equally.
	
	D
	Whether the government should support research with less immediate value.
	
	2.
	
	A
	Such research is also aimed at discovering truths.
	
	B
	Much of such research has proven beneficial to society.
	
	C
	Such research can have a long-term impact on our society.
	
	D
	Failure to support such research hinders the discovery of new knowledge.
	
	3.
	
	A
	Computer research was given priority even in the face of opposition.
	
	B
	Computer research was not well-received at first but has proven beneficial.
	
	C
	People have always been aware of the practical benefits of computer research.
	
	D
	The progress of computer research was slow at first due to insufficient funding.
\end{enumerate}

\subsubsection{Lectures 1}

\begin{enumerate}
	DBC
	1. Why is daydreaming often criticized?
	2. What is the benefit of daydreaming, according to the lecture?
	3. Why is it necessary to exercise control over daydreaming?
	1.
	
	A
	People often get lost in daydreams.
	
	B
	Too many people indulge in daydreaming.
	
	C
	Daydreaming indicates a lack of productivity.
	
	D
	People are taught that focus is the key to success.
	
	2.
	
	A
	It makes our brains more active.
	
	B
	It helps us generate original ideas.
	
	C
	It creates endless possibilities for success.
	
	D
	It helps us understand the people around us.
	
	3.
	
	A
	Daydreaming can waste time if it occurs too frequently.
	
	B
	Daydreaming may lead to a loss of one's self-awareness.
	
	C
	Daydreaming can distract one's attention from the current task.
	
	D
	Daydreaming may foster a constant desire to escape from the present.
\end{enumerate}


\subsubsection{Lectures 2}

\begin{enumerate}
	AADB
	1. What can we learn about being artistic from the lecture?
	2. What is creativity, according to the lecture?
	3. What does the speaker say about innovation?
	4. What is the difference between creativity and innovation?
	1.
	
	A
	It involves developing skills and talents.
	
	B
	It is the ability to think outside the box.
	
	C
	It is about doing something that's unknown to others.
	
	D
	It means creating for specific listeners, viewers, or users.
	
	2.
	
	A
	It involves the development of new ideas.
	
	B
	It involves self-reflection or self-expression.
	
	C
	It is about turning imaginative ideas into reality.
	
	D
	It is about employing artistic skills to create something.
	
	3.
	
	A
	It is more of an imaginative process than a productive one.
	
	B
	It represents a spontaneous departure from the mainstream.
	
	C
	It often leads to outcomes that can be seen in various forms.
	
	D
	It is about creating something new that can bring value to others.
	
	4.
	
	A
	Innovation is a process; creativity is the final product.
	
	B
	Creativity helps generate ideas; innovation applies them.
	
	C
	Innovation is linked to science; creativity is linked to art.
	
	D
	Creativity is inward-focused; innovation is outward-focused.
\end{enumerate}

\newpage
\subsection{Unit 5}

\subsubsection{Conversation}

\begin{enumerate}
	BBDC
	1. Why has the man packed up?
	2. What is the man interested in?
	3. What does the man plan to do on weekends while in France?
	4. What does the woman decide to do, according to the conversation?
	1.
	
	A
	He is going to visit a friend.
	
	B
	He is going on a working holiday.
	
	C
	He is going on a business trip to France.
	
	D
	He has accepted a full-time job on a farm.
	
	2.
	
	A
	Doing housework.
	
	B
	Doing building work.
	
	C
	Working in the fields.
	
	D
	Working in the garden.
	
	3.
	
	A
	Stay in the local town.
	
	B
	Do extra work on the farm.
	
	C
	Go to Paris every weekend.
	
	D
	Visit a few places in France.
	
	4.
	
	A
	Find a place that interests her.
	
	B
	Advise the man to give up his plan.
	
	C
	Join the man for the unique experience.
	
	D
	Search the Internet for travel information.
\end{enumerate}


\subsubsection{Passage}

\begin{enumerate}
	ADC
	1. Why did the speaker become dissatisfied with her job at the company?
	2. What can help us achieve a sense of fulfillment, according to the speaker?
	3. What job does the speaker find most suitable for herself now?
	1.
	
	A
	She no longer felt passionate.
	
	B
	There was no salary increase.
	
	C
	She felt there was no flexibility.
	
	D
	Her boss was not understanding.
	
	2.
	
	A
	Starting our own businesses.
	
	B
	Creating something from scratch.
	
	C
	Turning hobbies into professional pursuits.
	
	D
	Pursuing anything that reflects our passion.
	
	3.
	
	A
	A business owner.
	
	B
	An online teacher.
	
	C
	A freelance writer.
	
	D
	An investment advisor.
\end{enumerate}

\subsubsection{Lectures 1}

\begin{enumerate}
	CCDD
	1. Why did the speaker find it rewarding to interpret for the Spanish author?
	2. Why can being an interpreter broaden one's horizons?
	3. Why do interpreters have the privilege of witnessing history unfold before their eyes?
	4. How does interpreting contribute to one's personal development?
	1.
	
	A
	He had a crucial role in the conference.
	
	B
	He was praised by the author for his good work.
	
	C
	He appreciated the audience's reactions to his work.
	
	D
	He accurately interpreted the author's ideas and works.
	
	2.
	
	A
	Interpreters are on a never-ending journey to different places.
	
	B
	Interpreters play different roles in various events and conferences.
	
	C
	Interpreters can meet people from diverse backgrounds and cultures.
	
	D
	Interpreters can acquire world views distinct from those of other people.
	
	3.
	
	A
	They help organize many high-level conferences.
	
	B
	They help make critical decisions in important negotiations.
	
	C
	They often have the best seats at many high-level conferences.
	
	D
	They enable effective communication during conferences and negotiations.
	
	4.
	
	A
	It makes one understand others better.
	
	B
	It helps one manage emotions effectively.
	
	C
	It helps one develop various speaking styles.
	
	D
	It helps one maintain composure under pressure.
\end{enumerate}


\subsubsection{Lectures 2}

\begin{enumerate}
	BDB
	1. What can we learn about the "in-between" space from the lecture?
	2. According to the lecture, what may undermine the positive effects of commuting?
	3. Which of the following is mentioned in the lecture as a means of making our commute more positive and fulfilling?
	1.
	
	A
	It leads to increased levels of stress.
	
	B
	It provides time for detachment and recovery.
	
	C
	It blurs the boundaries between work and home life.
	
	D
	It encourages greater mental engagement and productivity.
	
	2.
	
	A
	Focusing on the road map.
	
	B
	Getting off at the wrong bus stop.
	
	C
	Taking a longer commute every day.
	
	D
	Dwelling on the negatives of the workday.
	
	3.
	
	A
	Spending less time on the road.
	
	B
	Engaging in conversations with friends.
	
	C
	Watching the changing scenery along the way.
	
	D
	Choosing comfortable means of transportation.
\end{enumerate}


\newpage
\subsection{Unit 6}

\subsubsection{Conversation}

\begin{enumerate}
	ABCD
	1. What is the man's problem?
	2. Why doesn't the man want to ask Dr. Smith for help?
	3. What can we learn about the teaching assistant from the conversation?
	4. What will the man probably do, according to the conversation?
	1.
	
	A
	He is struggling with a test.
	
	B
	He hasn't finished reading his textbooks.
	
	C
	He has trouble writing psychology essays.
	
	D
	He thinks there are too many concepts in the textbook.
	
	2.
	
	A
	He is afraid of Dr. Smith.
	
	B
	He is too shy and nervous.
	
	C
	He failed Dr. Smith's class once before.
	
	D
	He is often called on by Dr. Smith in class.
	
	3.
	
	A
	She is a great lecturer.
	
	B
	She has office hours on Wednesdays.
	
	C
	She is good at making complex concepts easy.
	
	D
	She is known for being willing to help students.
	
	4.
	
	A
	Post his difficulties online for suggestions.
	
	B
	Do some research to prepare for the meeting with Dr. Smith.
	
	C
	Talk to some professors who have posted video lectures online.
	
	D
	Speak with the teaching assistant and watch some online lectures.
\end{enumerate}


\subsubsection{Passage}

\begin{enumerate}
	DCDB
	1. Why is it important to consider problems from other people's perspectives?
	2. What is a major cause of conflicts in relationships, according to the passage?
	3. How should we solve problems through communication?
	4. What is the passage mainly about?
	1.
	
	A
	It helps us agree with other people's points.
	
	B
	It helps us be better understood by other people.
	
	C
	It helps us understand the backgrounds of other people.
	
	D
	It helps us understand the reasons behind other people's points.
	
	2.
	
	A
	We don't respect other people's decisions.
	
	B
	We get upset with other people's misbehavior.
	
	C
	We expect other people to behave in a certain way.
	
	D
	We blame other people for their mistakes and limitations.
	
	3.
	
	A
	By looking for things that are left unsaid.
	
	B
	By bringing up old conflicts and solving them.
	
	C
	By talking frankly and honestly to earn each other's trust.
	
	D
	By focusing on positive issues and seeking common ground.
	
	4.
	
	A
	Different types of conflicts in our relationships.
	
	B
	Suggestions on dealing with conflicts in our relationships.
	
	C
	Suggestions on how to avoid conflicts in our relationships.
	
	D
	Reasons behind various types of conflicts in our relationships.
\end{enumerate}

\subsubsection{Lectures 1}

\begin{enumerate}
	CCB
	1. Why is styrofoam an environmental problem?
	2. How did the team feel at the beginning of the experiment?
	3. What is the team's solution for the problem of used styrofoam?
	1.
	
	A
	It ends up as litter.
	
	B
	It is impossible to degrade.
	
	C
	It is slow to degrade and hard to recycle.
	
	D
	Its manufacturing process is harmful to the environment.
	
	2.
	
	A
	They were encouraged by their immediate success.
	
	B
	They were doubtful about the method of their experiment.
	
	C
	They were discouraged by the numerous failures they faced.
	
	D
	They were relieved as they found the process easier than expected.
	
	3.
	
	A
	Heating it to make it vaporize.
	
	B
	Creating activated carbon from it.
	
	C
	Using several methods to recycle it.
	
	D
	Finding a use for it in manufacturing other products.
\end{enumerate}


\subsubsection{Lectures 2}

\begin{enumerate}
	CCB
	1. Why is styrofoam an environmental problem?
	2. How did the team feel at the beginning of the experiment?
	3. What is the team's solution for the problem of used styrofoam?
	1.
	
	A
	They focus on the relationships between objects.
	
	B
	They are good at making well-informed choices.
	
	C
	They place a high value on empathy and creativity.
	
	D
	They focus on individual objects and their characteristics.
	
	2.
	
	A
	Finance.
	
	B
	Science.
	
	C
	Psychology.
	
	D
	Engineering.
	
	3.
	
	A
	Logical reasoning.
	
	B
	Harmony and balance.
	
	C
	The acceptance of changes.
	
	D
	The adaptability of thinking patterns.
	
	4.
	
	A
	People are born either analytic or holistic thinkers.
	
	B
	People acquire their thinking styles from the environment.
	
	C
	People's thinking styles can be influenced by personal preferences.
	
	D
	People develop their thinking styles by learning from their experiences.
\end{enumerate}
